
% https://akaedu.github.io/book/ch13.html

\chapter{本阶段总结}
\BookAnchor{db.ch13}\BookAnchor{db.ch13.id2752091}%
善于学习的人都应该善于总结。本书的编排顺序充分考虑到知识的前后依赖关系，保证在讲解每个新知识点的时候都只用到前面章节讲过的知识，但正因为如此，很多相互关联的知识点被拆散到多个章节中了。我们一章一章地纵向学习过来之后，应该理出几个横切面，把拆散到各章节中的知识点串起来。请从以下几个方面整理和复习。



1、C的语法规则。


\begin{enumerate}

\item 源文件中所有函数定义之外可以出现哪些语法元素？

\item 函数定义之中可以出现哪些语法元素？

\item 语句有哪几种？

\item 哪些语法元素需要遵循标识符的命名规则？

\item 表达式由哪些语法元素组成？

\item 到目前为止学过哪些运算符？它们的优先级和结合性是怎样的？

\item 哪些运算符取操作数的左值？哪些运算符的操作数必须是整型？哪些运算符有Side Effect？

\item 哪些表达式可以做左值？哪些表达式只能做右值？

\item 哪些地方必须用常量表达式？哪些地方必须用整数常量表达式？
\end{enumerate}


2、思维方法与编程思想。


\begin{itemize}

\item 以概念为中心，\hyperlink{db.intro.program}{第 1 节 “程序和编程语言”}

\item 组合规则，\hyperlink{db.expr.expression}{第 5 节 “表达式”}

\item Least Surprise，\hyperlink{db.ch03s03.id2713785}{第 3 节 “形参和实参”}

\item 充分条件与必要条件，\hyperlink{db.ch03s04.id2714471}{第 4 节 “全局变量、局部变量和作用域”}

\item 封装，\hyperlink{db.ch04s02.id2718511}{第 2 节 “if/\BookSplit{}else语句”}

\item 布尔逻辑，\hyperlink{db.ch04s03.id2718970}{第 3 节 “布尔代数”}

\item 递归，\hyperlink{db.ch05s03.id2722951}{第 3 节 “递归”}

\item 函数式编程，\hyperlink{db.ch06s01.id2726104}{第 1 节 “while语句”}

\item 迭代（\hyperlink{db.ch06.id2726092}{第 6 章 \emph{循环语句}}）与增量式求解（\hyperlink{db.ch11s02.id2745253}{第 2 节 “插入排序”}）

\item 抽象，\hyperlink{db.ch07s02.id2730921}{第 2 节 “数据抽象”}

\item 数据驱动，\hyperlink{db.ch08s05.id2734861}{第 5 节 “多维数组”}

\item 分而治之，\hyperlink{db.ch11s04.id2746095}{第 4 节 “归并排序”}

\item 折半查找，\hyperlink{db.ch11s06.id2746698}{第 6 节 “折半查找”}

\item 回溯，\hyperlink{db.ch12s03.id2750166}{例 12.\BookSplit{}3 “用深度优先搜索解迷宫问题”}
\end{itemize}


3、调试方法


\begin{itemize}

\item 编译错误、运行时错误与语义错误，\hyperlink{db.ch01s03.id2704937}{第 3 节 “程序的调试”}

\item 增量式开发，\hyperlink{db.ch05s02.id2721551}{第 2 节 “增量式开发”}

\item 打印语句与Scaffold，\hyperlink{db.ch05s02.id2721551}{第 2 节 “增量式开发”}

\item gdb，\hyperlink{db.ch10.id2740153}{第 10 章 \emph{gdb}}

\item DbC与Assertion，\hyperlink{db.ch11s06.id2746698}{第 6 节 “折半查找”}
\end{itemize}

\BookNav{\hyperlink{db.ch12s05}{上一页}\enspace\textperiodcentered\enspace \hyperlink{db.pt01}{上一级}\enspace\textperiodcentered\enspace \hyperlink{db.index}{起始页}\enspace\textperiodcentered\enspace \hyperlink{db.pt02}{下一页}}
